% Consolidated 38-radioactivity from 3014fceb998eba4990e2b4a247494872c4c8e1cc; original TeX fragments preserved.

% === BEGIN 38-radioactivity-214-226.tex blob ec7b89836441f9abdfd87dadaa1822a2bc827c2d ===
% 原书第六篇6.5节：放射性214--226；印刷页257，PDF266。
% 底本书末答案印刷348--349，PDF357--358。
\radioactivityproblem{214}{257}
放射性核衰变常数为 $\lambda$。
求 (a) 由 $t=0$ 到 $t$ 衰变的概率；
(b) 核的平均寿命。
\exsolution
单核在 $[0,t]$ 内不衰变的生存概率
$S(t)=e^{-\lambda t}$，
因此
\[
 \boxed{P(0<T_d<t)=1-e^{-\lambda t}}.
\]
衰变时间密度
$f(t)=\lambda e^{-\lambda t}$，故
\[
 \boxed{\tau=\int_0^\infty t f(t)\,dt
 =\int_0^\infty e^{-\lambda t}dt
 =\frac1{\lambda}}.
\]
\excheck 概率极限 $P(0)=0$、
$P(\infty)=1$；均寿命不等于半衰期。

\radioactivityproblem{215}{257}
半衰期 $T_{1/2}=71.3\,\mathrm d$ 的
放射性钴原子核，一个月内有多少份额发生衰变？
\exsolution
衰变常数
$\lambda=\ln2/T_{1/2}$。
按一个月约30天，
剩余份额 $2^{-30/71.3}$，
所以
\[
 \boxed{f=1-2^{-30/71.3}\simeq0.253
 \approx\frac14}.
\]
\excheck “一个月”并非固定秒数，
采用30天与实际公历月的近似差小于本题精度。

\radioactivityproblem{216}{257}
质量 $m=1.0\,\mu\mathrm g$ 的
$\mathrm{^{24}Na}$，半衰期15小时，
一小时内发射多少粒 $\beta$ 粒子？
\exsolution
初始核数
$N_0=(m/M)N_A$，
$M\simeq24\,\mathrm{g/mol}$。
假定每次 $\beta$ 衰变给出一电子，
一小时内衰变数
\[
 \boxed{\Delta N=
 \frac{mN_A}{M}
 \left(1-2^{-1/15}\right)
 \simeq1.13\times10^{15}\approx1.2\times10^{15}}.
\]
\excheck 不可用 $A_0\,t$ 当作精确值，
因为一小时内活度本身也随时间减少。

\radioactivityproblem{217}{257}
从 $t=0$ 开始记录
$\mathrm{^{23}Mg}$ 的 $\beta$ 衰变。
$t_1=2.0\,\mathrm s$ 内计得 $N_1$ 粒；
到 $t_2=3t_1$ 时累计计数
$N_2=2.66N_1$。
求平均寿命。
\exsolution
$N(t)=N_0(1-e^{-\lambda t})$。
令 $x=e^{-\lambda t_1}$，
由于 $t_2=3t_1$，
\[
 2.66=\frac{1-x^3}{1-x}=1+x+x^2.
\]
取 $0<x<1$ 的根
$x=[-1+\sqrt{1+4(2.66-1)}]/2\simeq0.882$。
故
\[
 \boxed{\tau=\frac1\lambda=
 -\frac{t_1}{\ln x}\simeq15.9\,\mathrm s
 \approx16\,\mathrm s}.
\]
\excheck 此题是\emph{累计计数}的比例，
不是瞬时活度的比例。

\radioactivityproblem{218}{257}
某放射性制剂的活度经过7.0天降低为原来的
$1/2.5$。求半衰期。
\exsolution
$A(t)/A(0)=e^{-\lambda t}=1/2.5$，
所以 $\lambda=\ln2.5/7.0\,\mathrm{d^{-1}}$。
由 $T_{1/2}=\ln2/\lambda$，
\[
 \boxed{T_{1/2}=
 \frac{7.0\,\mathrm d\,\ln2}{\ln2.5}
 \simeq5.30\,\mathrm d}.
\]
\excheck 活度与未衰变核数均随同一指数衰减。

\radioactivityproblem{219}{257}
某放射性核素初始活度为
$A_0=650\,\mathrm{min^{-1}}$。
经过\emph{半个半衰期}后活度是多少？
\exsolution
$A(t)=A_02^{-t/T_{1/2}}$。
当 $t=T_{1/2}/2$，
\[
 \boxed{A=650\,2^{-1/2}
 \simeq4.60\times10^2\,\mathrm{min^{-1}}}.
\]
\excheck 经过半个半衰期不是活度减半；
应乘 $2^{-1/2}$。

\radioactivityproblem{220}{257}
某 $\mathrm{^{55}Co}$ 制剂的活度每小时下降4.0\%，
产物不再有放射性。求衰变常数与平均寿命。
\exsolution
一小时后活度比
$A(1\,\mathrm h)/A_0=0.960$，
故
\[
 \boxed{\lambda=-\frac{\ln0.960}{3600\,\mathrm s}
 \simeq1.134\times10^{-5}\,\mathrm{s^{-1}}}.
\]
由 $\tau=1/\lambda$，
\[
 \boxed{\tau\simeq8.82\times10^4\,\mathrm s
 =24.5\,\mathrm h\simeq1.02\,\mathrm d}.
\]
\excheck 本版书末将均寿命写成“1.0 years”；
照印刷的\emph{每小时下降4\%}应得约1.0天，
而不可能是一年。题目参数与书末单位冲突，已记录。

\radioactivityproblem{221}{257}
质量1.0g的$\mathrm{^{238}U}$制剂
每秒发射 $1.24\times10^4$ 粒 $\alpha$ 粒子。
求半衰期及制剂活度。
\exsolution
一克纯 $\mathrm{^{238}U}$ 核数
$N=(1\,\mathrm g/238\,\mathrm{g/mol})N_A$。
若每次衰变发射一粒 $\alpha$，
$A=\lambda N$，所以
\[
 \boxed{T_{1/2}=\frac{N\ln2}{A}
 \simeq4.49\times10^9\,\mathrm{year}}.
\]
活度即给出的衰变率
\[
 \boxed{A=1.24\times10^4\,\mathrm{Bq}}.
\]
\excheck $\mathrm{Bq}=1\,\mathrm{s^{-1}}$；
不可把 $\alpha$ 计数与单位换算再乘一分钟。

\radioactivityproblem{222}{257}
古木中 $\mathrm{^{14}C}$ 比活度为新砍伐木材
比活度的 $3/5$，
已知 $\mathrm{^{14}C}$ 半衰期为5570年，
估算古木年代。
\exsolution
假定碳交换停止之后 $\mathrm{^{14}C}$ 只衰变，
\[
 \frac{A(t)}{A(0)}=e^{-\lambda t}=\frac35,\quad
 \lambda=\frac{\ln2}{5570\,\mathrm{yr}}.
\]
故
\[
 \boxed{t=\frac{5570}{\ln2}\ln\frac53
 \simeq4.11\times10^3\,\mathrm{yr}}.
\]
\excheck 需要假设最初与同环境新木的
$\mathrm{^{14}C}$ 比活度相同；
否则须进行年代校准。

\radioactivityproblem{223}{257}
某铀矿中 $\mathrm{^{238}U}$ 核数与
$\mathrm{^{206}Pb}$ 核数之比为
$\eta=2.8$。
假定矿中所有 $\mathrm{^{206}Pb}$
均来自该铀衰变链，
$T_{1/2}=4.5\times10^9$ 年，
估算铀矿年龄。
\exsolution
无初始铅且铀衰变链子代相对迅速达到稳定铅时，
$N_U=N_0e^{-\lambda t}$、
$N_{\rm Pb}=N_0-N_U$。
由 $N_U/N_{\rm Pb}=\eta$ 得
$e^{\lambda t}=1+1/\eta$，因此
\[
 \boxed{t=\frac{T_{1/2}}{\ln2}
 \ln\left(1+\frac1\eta\right)
 \simeq2.00\times10^9\,\mathrm{yr}}.
\]
\excheck 必须假定初始没有额外稳定铅，
否则不能仅凭测量的核数比唯一确定矿龄。

\radioactivityproblem{224}{257}
计算 $\mathrm{^{24}Na}$ 与 $\mathrm{^{235}U}$
的比活度，已知半衰期分别为15小时与
$7.1\times10^8$ 年。
\exsolution
一克纯放射性核素中
$N=N_A/M$，而 $A=\lambda N$。
故比活度
$a_s=(\ln2/T_{1/2})(N_A/M)$。
代入 $M_{\rm Na}=24\,\mathrm{g/mol}$、
$M_{\rm U}=235\,\mathrm{g/mol}$，
\[
 \boxed{a_s(\mathrm{^{24}Na})\simeq
 3.2\times10^{17}\,\mathrm{Bq/g}},
\quad
 \boxed{a_s(\mathrm{^{235}U})\simeq
 8.0\times10^{4}\,\mathrm{Bq/g}}.
\]
\excheck 比活度单位为单位质量每秒衰变数；
年代须先换成秒。

\radioactivityproblem{225}{257}
把活度 $A_0=2.0\times10^3\,\mathrm{Bq}$
的 $\mathrm{^{24}Na}$ 溶液注入血液。
$t=5.0$小时后，
每 $\mathrm{cm^3}$ 血液样本活度为
$A'=16\,\mathrm{min^{-1}\,cm^{-3}}$。
该核素半衰期15小时，求血液总体积。
\exsolution
假定核素在血液中均匀分布且无体外排泄，
$t$时剩余总活度为
$A_0\,2^{-t/T_{1/2}}$。
每立方厘米样本活度
$A'=16/60\,\mathrm{Bq/cm^3}$，
所以
\[
 \boxed{V=\frac{A_0}{A'}\,2^{-t/T_{1/2}}
 =\frac{2000}{16/60}\,2^{-5/15}
 \simeq5.95\times10^3\,\mathrm{cm^3}
 \approx6.0\,\mathrm L}.
\]
\excheck 实际医学使用须考虑组织分布、
排泄及剂量安全；这里纯属理想示踪守恒模型。

\radioactivityproblem{226}{257}
含放射性 $\mathrm{^{58}Co}$ 和非放射性
$\mathrm{^{59}Co}$ 的制剂，
整体比活度为
$2.2\times10^{12}\,\mathrm{Bq/g}$。
$\mathrm{^{58}Co}$ 半衰期为71.3天，
求放射性钴所占质量百分数。
\exsolution
若每克纯 $\mathrm{^{58}Co}$ 的比活度
$a_{58}=(\ln2/T_{1/2})(N_A/58)$，
混合物的整体比活度为 $a=f a_{58}$，
其中 $f$ 为放射性同位素质量份额。
因此
\[
 \boxed{f=\frac{a}{a_{58}}
 =\frac{aT_{1/2}\,58}{N_A\ln2}
 \simeq1.88\times10^{-3}},
 \quad \boxed{100f\simeq0.188\%\approx0.19\%}.
\]
\excheck 非放射性的 $\mathrm{^{59}Co}$
提供质量但不提供活度，所以直接按质量份额线性缩放。
% === END 38-radioactivity-214-226.tex ===

% === BEGIN 38-radioactivity-227-237.tex blob d8b00d0391eceedf1be3e56938b121fcc933c0c8 ===
% 原书第六篇6.5节227--237；印刷258--259，PDF267--268；
% 书末答案348--349，PDF357--358。
\radioactivityproblem{227}{258}
一制剂由两种半衰期不同的 $\beta$ 放射性核素组成。
按时间 $t$（小时）测得
\[
\begin{array}{c|rrrrrrrrr}
t&0&1&2&3&5&7&10&14&20\\\hline
\ln A&4.10&3.60&3.10&2.60&2.06&1.82&1.60&1.32&0.90
\end{array}
\]
估算两种核素的半衰期及 $t=0$ 时核数比。
\exsolution
总活度满足双指数模型
\[
 A(t)=A_{f0}e^{-\lambda_f t}
 +A_{s0}e^{-\lambda_s t},\quad\lambda_f>\lambda_s.
\]
末期 $t=14,20$ 两个点给慢分量的近似斜率
$\lambda_s\simeq(1.32-0.90)/6=0.070\,\mathrm{h^{-1}}$，
于是 $A_{s0}\simeq e^{0.90+20\lambda_s}\simeq10$。
初始总活度为 $e^{4.10}\simeq60.34$，
所以 $A_{f0}\simeq50.34$。
从早期点减去慢分量再拟合快指数，得
$\lambda_f\simeq0.6$--$0.7\,\mathrm{h^{-1}}$。
若对全部九个 $\ln A$ 数据作双指数最小二乘拟合，
一组近似参数为
$\lambda_f\simeq0.701\,\mathrm{h^{-1}}$、
$\lambda_s\simeq0.0640\,\mathrm{h^{-1}}$，
$A_{f0}\simeq53.8$、$A_{s0}\simeq9.01$。
于是
\[
 \boxed{T_f\simeq0.99\,\mathrm h,\qquad
 T_s\simeq10.8\,\mathrm h}.
\]
由于 $N=A/\lambda$，慢/快初始核数之比
\[
 \boxed{\frac{N_{s0}}{N_{f0}}=
 \frac{A_{s0}\lambda_f}{A_{f0}\lambda_s}
 \simeq1.83}.
\]
\excheck 书末印 $T_f=1.6\,\mathrm h$、
$T_s=9.8\,\mathrm h$、核数比10。
把表中九点均按自然对数处理，
当前双指数参数的对数残差均方根约为
$\mathrm{RMS}_{\rm fitted}=0.0337$。
若保持书末三项参数不变，只允许整体活度
随$t=0$观测值归一化，则九点对数残差均方根
$\mathrm{RMS}_{\rm book}\simeq1.067$，
明显不符合原印数据。
这里的0.0337也高于仅有0.01量级的表格
末位四舍五入误差，说明“拟合值”并非精确通过
所有观测点；不能误标为无误差的代数反演。
本稿给出的是数值最小二乘近似，
保留原书答案与原始表格的冲突供二校。

\radioactivityproblem{228}{258}
反应堆以恒定速率
$q=2.7\times10^9\,\mathrm{s^{-1}}$
生产 $\mathrm{^{32}P}$，
半衰期 $T_{1/2}=14.3\,\mathrm d$。
开始生产后多久该核素的活度达到
$A=1.0\times10^9\,\mathrm{Bq}$？
\exsolution
由生产与衰变的核数方程
$dN/dt=q-\lambda N$，$N(0)=0$，
解得
$N(t)=(q/\lambda)(1-e^{-\lambda t})$，
活度
$A(t)=\lambda N(t)=q(1-e^{-\lambda t})$。
故
\[
 \boxed{t=-\frac1\lambda\ln\left(1-\frac Aq\right)
 =-\frac{T_{1/2}}{\ln2}\ln\left(1-\frac Aq\right)
 \simeq9.5\,\mathrm d}.
\]
\excheck 随时间 $A(t)\to q$，
目标活度必须满足 $A<q$，本题成立。

\radioactivityproblem{229}{258}
一条衰变链 $A_1\xrightarrow{\lambda_1}A_2
\xrightarrow{\lambda_2}\cdots$，
初始只有 $N_{10}$ 个母核 $A_1$，
求 (a) 子核 $A_2$ 随时间的数量；
(b) 子核活度达到最大时刻。
\exsolution
母核
$N_1=N_{10}e^{-\lambda_1t}$，
子核满足
$dN_2/dt=\lambda_1N_1-\lambda_2N_2$，
$N_2(0)=0$。
乘积分因子 $e^{\lambda_2t}$ 积分，得到
\[
 \boxed{N_2(t)=\frac{\lambda_1N_{10}}
 {\lambda_2-\lambda_1}
 (e^{-\lambda_1t}-e^{-\lambda_2t})}.
\]
子核活度 $A_2=\lambda_2N_2$，在
$dN_2/dt=0$ 时极大；
$\lambda_1e^{-\lambda_1t_m}
 =\lambda_2e^{-\lambda_2t_m}$，
所以
\[
 \boxed{t_m=\frac{\ln(\lambda_1/\lambda_2)}
 {\lambda_1-\lambda_2}}.
\]
\excheck 对 $\lambda_1\ne\lambda_2$ 公式成立，
并且 $t_m>0$；两常数相等时取极限见下题。

\radioactivityproblem{230}{258}
若前题两衰变常数相等
$\lambda_1=\lambda_2=\lambda$，
求子核 $A_2$ 数量随时间的函数及达到最大值的时刻。
\exsolution
方程
$dN_2/dt+\lambda N_2=
\lambda N_{10}e^{-\lambda t}$，
乘 $e^{\lambda t}$ 得
$d(N_2e^{\lambda t})/dt=\lambda N_{10}$。
因此
\[
 \boxed{N_2(t)=\lambda N_{10}t e^{-\lambda t}}.
\]
求导
$dN_2/dt=\lambda N_{10}e^{-\lambda t}(1-\lambda t)$，
所以
$\boxed{t_m=1/\lambda=\tau}$。
\excheck 两指数相减型
Bateman公式在 $\lambda_2\to\lambda_1$
时也趋于此式。

\radioactivityproblem{231}{258}
衰变链
$A_1\xrightarrow{\lambda_1}A_2
\xrightarrow{\lambda_2}A_3$
中 $A_3$ 稳定，初始只有 $N_{10}$ 个 $A_1$。
求稳定核 $A_3$ 的累积数目。
\exsolution
根据上题组解，
$N_1=N_{10}e^{-\lambda_1 t}$、
\[
 N_2=\frac{\lambda_1N_{10}}
 {\lambda_2-\lambda_1}
 (e^{-\lambda_1t}-e^{-\lambda_2t}).
\]
整个核链核数守恒，
$N_1+N_2+N_3=N_{10}$，
所以
\[
 \boxed{N_3(t)=N_{10}\left[
 1+\frac{\lambda_1e^{-\lambda_2t}
 -\lambda_2e^{-\lambda_1t}}
 {\lambda_2-\lambda_1}\right]}.
\]
\excheck $N_3(0)=0$，而
$\lim_{t\to\infty}N_3(t)=N_{10}$，
符合所有母核最终变为稳定核的假定。

\radioactivityproblem{232}{258}
$\mathrm{^{210}Bi}\xrightarrow[\beta^-]{\lambda_1}
\mathrm{^{210}Po}\xrightarrow[\alpha]{\lambda_2}
\mathrm{^{206}Pb}$（稳定），
$\lambda_1=1.60\times10^{-6}\,\mathrm{s^{-1}}$、
$\lambda_2=5.80\times10^{-8}\,\mathrm{s^{-1}}$。
初始纯 $\mathrm{^{210}Bi}$ 的质量 $1.00\,\mathrm{mg}$。
求制成一个月后的 $\alpha$、$\beta$ 活度。
\exsolution
初始母核数
$N_0=(10^{-3}\,\mathrm g/210\,\mathrm{g/mol})N_A$。
取一月30天，$t=30(86400)\,\mathrm s$。
母核数 $N_{\rm Bi}=N_0e^{-\lambda_1t}$，
因此
\[
 \boxed{A_\beta=\lambda_1N_0e^{-\lambda_1t}
 \simeq7.25\times10^{10}\,\mathrm{Bq}}.
\]
子核
\[
 N_{\rm Po}=
 \frac{\lambda_1N_0}{\lambda_2-\lambda_1}
 (e^{-\lambda_1t}-e^{-\lambda_2t}),
\]
所以
\[
 \boxed{A_\alpha=\lambda_2 N_{\rm Po}
 \simeq1.46\times10^{11}\,\mathrm{Bq}}.
\]
\excheck 题目给出两种不同衰变的活度，
应分别使用瞬时母、子核数，
而不是将初始母核数直接乘两个常数。

\radioactivityproblem{233}{258}
(a) $\mathrm{^{226}Ra}$ 经五次 $\alpha$
和四次 $\beta^-$ 衰变生成何核？
(b) $\mathrm{^{238}U}$ 最终变为稳定
$\mathrm{^{206}Pb}$，经历多少次 $\alpha$
与 $\beta^-$ 衰变？
\exsolution
一次 $\alpha$ 衰变使
$(A,Z)\mapsto(A-4,Z-2)$；
一次 $\beta^-$ 使
$(A,Z)\mapsto(A,Z+1)$。
(a) Ra 初始 $(226,88)$，
五次 $\alpha$ 后 $(206,78)$，
四次 $\beta^-$ 后 $(206,82)$，
故 $\boxed{\mathrm{^{206}Pb}}$。
(b) 质量数从238到206，减少32，
需要 $\boxed{n_\alpha=32/4=8}$。
铀初始 $Z=92$，
八次 $\alpha$ 后 $Z=76$，
到铅$Z=82$还需
$\boxed{n_{\beta^-}=6}$。
\excheck $\beta^-$ 不改变质量数，
因此先由质量数差求 $\alpha$ 次数更直接。

\radioactivityproblem{234}{258}
假设初始静止的 $\mathrm{^{206}Pb}$ 核
发生 $\alpha$ 衰变，粒子动能
$T_\alpha=5.77\,\mathrm{MeV}$。
求剩余娘核的反冲速度，
以及反冲能占释放总能量的份额。
\exsolution
衰变前总动量为零，因此娘核动量与 $\alpha$
动量大小相等反向。
非相对论近似
$p_\alpha=\sqrt{2m_\alpha T_\alpha}$，
故娘核速率
\[
 \boxed{v_d=\frac{\sqrt{2m_\alpha T_\alpha}}{M_d}
 \simeq3.3\times10^5\,\mathrm{m/s}},
 \quad M_d\simeq202u.
\]
反冲能比
$T_d/T_\alpha=m_\alpha/M_d$，
故在总动能
$Q=T_\alpha+T_d$ 中，
\[
 \boxed{\frac{T_d}{Q}
 =\frac{m_\alpha}{M_d+m_\alpha}
 \simeq\frac4{206}=1.94\%\approx2.0\%}.
\]
\excheck 本题按原书是\emph{假想}的
$\mathrm{^{206}Pb}$ 衰变运动学模型，
并不意味着稳定 $\mathrm{^{206}Pb}$ 真实自发 $\alpha$ 衰变。

\radioactivityproblem{235}{258--259}
质量 $1.00\,\mathrm{mg}$ 的 $\mathrm{^{210}Po}$
制剂在其平均寿命时间内释放多少热？
每次 $\alpha$ 粒子动能为 $5.3\,\mathrm{MeV}$，
娘核几乎都形成基态。
\exsolution
平均寿命 $\tau=1/\lambda$ 时，
已衰变份额
$f=1-e^{-1}$。
核数
$N_0=(m/M)N_A$。
因娘核反冲，
总释放能量略大于 $\alpha$ 动能：
$Q=T_\alpha(1+m_\alpha/M_d)
\simeq T_\alpha(210/206)$。
于是
\[
 \boxed{W=N_0(1-e^{-1})Q
 \simeq1.57\times10^6\,\mathrm J
 \approx1.6\,\mathrm{MJ}}.
\]
\excheck 经过一个平均寿命仅约63.2\%的核衰变，
不能把所有核都当作已经衰变。

\radioactivityproblem{236}{259}
$\mathrm{^{210}Po}$ 的 $\alpha$ 衰变有两组
$\alpha$ 动能 $5.30$ 与 $4.50\,\mathrm{MeV}$；
前者使娘核至基态，后者至激发态。
求退激时 $\gamma$ 光子的能量。
\exsolution
母、娘核静质量差相同，
两衰变分支的总释放能之差
等于娘核激发能。
反冲动能比为$m_\alpha/M_d$，
所以
\[
 E^*=(T_{\alpha,1}-T_{\alpha,2})
 \left(1+\frac{m_\alpha}{M_d}\right).
\]
取$M_d\simeq206u$，
\[
 \boxed{E_\gamma\simeq E^*
 =(5.30-4.50)\frac{210}{206}
 \simeq0.816\,\mathrm{MeV}\approx0.82\,\mathrm{MeV}}.
\]
\excheck $\gamma$ 发射也有微小核反冲，
本题 MeV 精度下可以忽略。

\radioactivityproblem{237}{259}
常温常压空气中 $\alpha$ 粒子平均射程经验式
$R=0.98\times10^{-27}v_0^3\,\mathrm{cm}$，
其中 $v_0$ 的单位为 $\mathrm{cm/s}$。
初始动能 $T_\alpha=7.0\,\mathrm{MeV}$。
求 (a)射程；(b)全程及前半段射程内
平均产生的离子对数，每对形成能为 $34\,\mathrm{eV}$。
\exsolution
初速度
$v_0=\sqrt{2T_\alpha/m_\alpha}
\simeq1.84\times10^9\,\mathrm{cm/s}$，
代经验式
\[
 \boxed{R\simeq6.1\,\mathrm{cm}}.
\]
全部动能转为电离的近似下，
\[
 \boxed{N_{\rm total}=\frac{7.0\times10^6}{34}
 \simeq2.06\times10^5}.
\]
为计算前半段，用残余射程
$R_{\rm rem}\propto v^3\propto T^{3/2}$，
因此通过前半段距离后残余射程为$R/2$，
剩余动能
$T_{\rm mid}=T_0\,2^{-2/3}$。
前半段的能量损失为
$\Delta T=T_0(1-2^{-2/3})$，所以
\[
 \boxed{N_{\rm first\,half}=
 \frac{7.0\times10^6}{34}(1-2^{-2/3})
 \simeq7.62\times10^4}.
\]
\excheck 能量损失沿射程并不均匀，
故前半段产生的离子对不是总数的一半。
% === END 38-radioactivity-227-237.tex ===

% === BEGIN 38-radioactivity-238-248.tex blob 56fe731462c4eae65822c2212194796abb952e33 ===
% 原书第六篇6.5节238--248（放射性本节完），印刷259--260，
% 用户原版PDF268--269，书末答案349，PDF358。
\radioactivityproblem{238}{259}
用母原子质量 $M_p$、娘原子质量 $M_d$
和电子质量 $m_e$，
推导 $\beta^-$、$\beta^+$ 与K俘获
放出的反应能 $Q$。
\exsolution
令原子质量包含其相应数目的束缚电子。
$\beta^-$ 衰变原子序数增加1，
核衰变伴随发射一电子，
母/娘中性原子的电子数差恰好抵消发射电子质量，
故
\[
 \boxed{Q_{\beta^-}=(M_p-M_d)c^2}.
\]
$\beta^+$ 衰变原子序数减1，同时发射正电子，
母原子与娘原子电子云计数以及正电子质量
共使静质量差需扣除 $2m_e$：
\[
 \boxed{Q_{\beta^+}=(M_p-M_d-2m_e)c^2}.
\]
K俘获时被俘获的电子已属于母原子，
原子质量的电子数恰能抵消，
\[
 \boxed{Q_{K}=(M_p-M_d)c^2}.
\]
\excheck 各式是中性\emph{原子质量}近似；
如研究俘获后特定壳层的中微子能量，还应扣除
电子空穴的激发能并考虑核反冲。

\radioactivityproblem{239}{259}
用原书原子质量表计算
$\mathrm{^{10}Be}$ 的 $\beta^-$ 衰变最大电子动能，
以及娘核直接处于基态时的反冲动能。
\exsolution
反应为
$\mathrm{^{10}Be}\to\mathrm{^{10}B}+e^-+\bar\nu_e$。
原版质量表中
$\Delta M_{\mathrm{^{10}Be}}\simeq0.01354u$，
$\Delta M_{\mathrm{^{10}B}}\simeq0.01294u$，
故
$Q\simeq(0.00060u)c^2\simeq0.559\,\mathrm{MeV}$。
当中微子能量趋于0时，
电子动能接近最大值，
\[
 \boxed{T_{e,\max}\simeq0.56\,\mathrm{MeV}}.
\]
此时娘核动量与电子相反；
相对论电子满足
$p_e^2c^2=T_e(T_e+2m_ec^2)$，
娘核非相对论反冲能
\[
 \boxed{T_R\simeq
 \frac{T_{e,\max}(T_{e,\max}+2m_ec^2)}
 {2M_dc^2}\simeq47.5\,\mathrm{eV}}.
\]
\excheck 计算反冲必须使用相对论电子动量，
不能简单取 $p_e^2=2m_eT_e$。

\radioactivityproblem{240}{259}
质量 $m=1.0\,\mathrm{mg}$ 的
$\mathrm{^{24}Na}$ 纯 $\beta^-$ 放射性制剂，
半衰期 $T_{1/2}=15$小时。
假设每个 $\beta$ 的平均动能为其端点能的1/3，
求一天内释放的热量。
\exsolution
反应
$\mathrm{^{24}Na}\to\mathrm{^{24}Mg}+e^-+\bar\nu$。
原书质量超额分别为
$-0.00903u$、$-0.01496u$，
所以
$Q=(0.00593u)c^2\simeq5.53\,\mathrm{MeV}$。
按题目假设
$\bar T_\beta=Q/3$。
初核数
$N_0=(10^{-3}\,\mathrm g/24\,\mathrm{g/mol})N_A$，
一天内衰变份额
$f=1-2^{-24/15}$。
故近似放热
\[
 \boxed{W=N_0 f\,\frac Q3
 \simeq5.0\times10^6\,\mathrm J}.
\]
\excheck 题目明确把 $\beta$ 平均动能指定为端点的1/3；
中微子逃逸能不计热量，
娘核反冲及可能随后的其他退激辐射亦忽略。

\radioactivityproblem{241}{259}
用原子质量表计算 $\mathrm{^{11}C}$ 的
$\beta^+$ 衰变中，当娘核\emph{不反冲}时
正电子与中微子的动能。
\exsolution
从原书表得质量超额
$\Delta_{\mathrm{^{11}C}}=0.01143u$、
$\Delta_{\mathrm{^{11}B}}=0.00930u$，
故正电子衰变可用动能
$Q=(0.00213u)c^2-2m_ec^2\simeq0.96\,\mathrm{MeV}$。
娘核零反冲要求正电子与中微子动量相等、方向相反。
令正电子动能 $T$，中微子能量 $E_\nu$，
则
\[
 T+E_\nu=Q,\quad
 E_\nu=pc=\sqrt{T(T+2m_ec^2)}.
\]
两式平方求解
\[
 \boxed{T=\frac{Q^2}{2(Q+m_ec^2)}
 \simeq0.32\,\mathrm{MeV}},
 \quad
 \boxed{E_\nu=Q-T\simeq0.65\,\mathrm{MeV}}.
\]
\excheck 由于正电子有静质量，而中微子近似无质量，
零反冲不意味着两者分得相同能量。

\radioactivityproblem{242}{259}
$\mathrm{^{13}N}$ 的 $\beta^+$ 衰变中，
当正电子动能最大时，
求娘核反冲动能。
\exsolution
衰变为 $\mathrm{^{13}N}\to
\mathrm{^{13}C}+e^++\nu_e$。
原书质量表质量超额
$0.00574u$ 与 $0.00335u$，
故
$Q=(0.00239u)c^2-2m_ec^2
 \simeq1.20\,\mathrm{MeV}$。
最大正电子动能对应中微子能量趋于零，
娘核动量 $p_R\simeq p_e$。
因此
\[
 \boxed{T_R\simeq
 \frac{Q(Q+2m_ec^2)}{2M_Rc^2}
 \simeq0.11\,\mathrm{keV}},
 \quad M_R\simeq13u.
\]
\excheck 这里核反冲能远小于 MeV，
可以在端点近似中令正电子动能 $T\simeq Q$。

\radioactivityproblem{243}{259}
使用原子质量表计算
$\mathrm{^{7}Be}$ 发生 K 电子俘获、
娘核处于基态时的反冲速度。
\exsolution
反应
$\mathrm{^{7}Be}+e^-\to\mathrm{^{7}Li}+\nu_e$。
原书质量超额分别
$0.01693u$ 与 $0.01601u$，
故 $Q\simeq(0.00092u)c^2
 \simeq0.857\,\mathrm{MeV}$。
中微子携带几乎全部释放能量，
动量为 $p_\nu\simeq Q/c$，
娘核反冲速率
\[
 \boxed{v_R\simeq\frac{Q}{M_Rc}
 \simeq3.9\times10^4\,\mathrm{m/s}
 \approx40\,\mathrm{km/s}}.
\]
\excheck K俘获后的壳层空穴可能影响
中微子精细谱线能量；本题要求的数量级可忽略。

\radioactivityproblem{244}{259}
激发 $\mathrm{^{109}Ag}$ 核退激时，
或发射能量 $87\,\mathrm{keV}$ 的 $\gamma$ 光子，
或发射K内转换电子，其结合能为
$26\,\mathrm{keV}$。
求这些电子的速率。
\exsolution
内转换过程将核跃迁能量直接交给壳层电子，
其动能近似
$T_e=87-26=61\,\mathrm{keV}$。
电子总能
$E_e=m_ec^2+T_e=\gamma m_ec^2$，
因此
\[
 \gamma=1+\frac{61}{511},\quad
 \boxed{v_e=c\sqrt{1-\gamma^{-2}}
 \simeq0.45c}.
\]
\excheck 内转换电子并非由 $\gamma$ 光子二次光电吸收产生，
而是核与原子电子直接耦合的退激通道。

\radioactivityproblem{245}{259}
一个初始静止的自由
$\mathrm{^{191}Ir}$ 激发核具有
$E^*=129\,\mathrm{keV}$，
向基态跃迁放出 $\gamma$ 光子。
核反冲引起光子能量的相对改变是多少？
\exsolution
若不计反冲光子能量为 $\varepsilon_0=E^*$。
发射后核动量 $p=\varepsilon/c$，
反冲能
$T_R\simeq\varepsilon^2/(2Mc^2)$。
能量守恒
$E^*=\varepsilon+T_R$。
因 $\varepsilon\ll Mc^2$，
相对光子能量降低为
\[
 \boxed{\frac{\varepsilon_0-\varepsilon}{\varepsilon_0}
 \simeq\frac{E^*}{2Mc^2}
 \simeq3.6\times10^{-7}}.
\]
\excheck 分母中是约 $191u\,c^2$ 的核静能，
不是 $129\,\mathrm{keV}$ 本身。

\radioactivityproblem{246}{259}
自由 $\mathrm{^{191}Ir}$ 核组成的光源与吸收体
发生 $129\,\mathrm{keV}$ 的共振吸收，
为抵消发射、吸收两端的核反冲能差，
两者应有多大相对速率？
\exsolution
自由核发射光子比核能级间隔少
$E_R=\varepsilon^2/(2Mc^2)$；
静止核吸收激发则要求比能级间隔多
$E_R$。两端光子能量失配约
$\Delta E=2E_R=\varepsilon^2/(Mc^2)$。
小速度多普勒能移
$|\Delta\varepsilon|\simeq\varepsilon v/c$，
令其补偿失配，
\[
 \boxed{v=\frac{\varepsilon}{Mc}
 \simeq2.2\times10^2\,\mathrm{m/s}
 =0.22\,\mathrm{km/s}}.
\]
\excheck 这里只求相对速率的模，
必须选择使吸收体观察到光子蓝移的相对运动方向。

\radioactivityproblem{247}{260}
$\gamma$ 射线源位于吸收体上方
$h=20\,\mathrm m$，
为了完全补偿光子向下传播时因地球重力
产生的能量变化，光源应以多大速率向上运动？
\exsolution
弱引力场、低速极限，
光子向下传播的引力能量增量相对为
$\Delta\varepsilon/\varepsilon\simeq gh/c^2$。
光源向上远离吸收体时，
其发射光子相对吸收体发生红移，
一阶相对量约 $v/c$。
令两个效应抵消，
$v/c=gh/c^2$，所以
\[
 \boxed{v=\frac{gh}{c}\simeq
 6.5\times10^{-7}\,\mathrm{m/s}
 =0.65\,\mu\mathrm{m/s}}.
\]
\excheck 问的是光源\emph{向上}运动，
用于抵消沿竖直向下传播光子的引力蓝移。

\radioactivityproblem{248}{260}
$\mathrm{^{67}Zn}$ 激发核发出能量
$\varepsilon=93\,\mathrm{keV}$ 的
Mössbauer 共振 $\gamma$ 光子，
激发态平均寿命 $\tau=14\,\mu\mathrm s$。
源应抬高至少多少，
使在地面观测到的引力谱线位移超过其自然线宽？
\exsolution
弱场引力能量相对变化
$\Delta\varepsilon/\varepsilon\simeq gh/c^2$，
绝对能移
$\Delta\varepsilon\simeq\varepsilon gh/c^2$。
激发态有限寿命导致自然线宽数量级
$\Gamma\simeq\hbar/\tau$。
令 $\Delta\varepsilon\gtrsim\Gamma$，
\[
 \boxed{h_{\min}\simeq
 \frac{\hbar c^2}{g\varepsilon\tau}}.
\]
取 $\varepsilon=93\,\mathrm{keV}$、
$\tau=14\times10^{-6}\,\mathrm s$，
得
$\boxed{h_{\min}\simeq4.6\,\mathrm m}$。
\excheck 使用 $\hbar$ 而非 $h$，
因为自然宽度与角频率寿命关系
$\Gamma\simeq\hbar/\tau$。
% === END 38-radioactivity-238-248.tex ===
